\answer{OP37}
Sort the points as $X_{(1)}<\cdots<X_{(N)}$ and define the interior gaps
\[
D_i=X_{(i+1)}-X_{(i)},\qquad 1\le i\le N-1.
\]
Mutual nearest neighbors must be adjacent. Let $I_i$ indicate that the endpoints of adjacent edge $i$ choose one another. The first edge qualifies precisely when $D_1<D_2$, and the last edge has the analogous condition. An interior edge, $2\le i\le N-2$, qualifies precisely when
\[
D_i<D_{i-1},\qquad D_i<D_{i+1}.
\]
All uniform spacings, including the two boundary spacings, are exchangeable and have zero probability of ties. Each end edge therefore qualifies with probability $1/2$, while an interior gap is smallest among the relevant three with probability $1/3$.

Each point has a unique nearest neighbor, so distinct mutual pairs cannot share an endpoint. Every pair contributes exactly two good points. Thus, if $G$ is the number of good points,
\[
G=2\sum_{i=1}^{N-1}I_i,
\quad
\E G=2\left(\frac12+\frac12+\frac{N-3}{3}\right)
=\boxed{\frac{2N}{3}}.
\]
For $N=3$ the interior sum is empty and the same formula applies. The expected number of pairs, rather than points, is $N/3$.
